% Gemeinsame Praeambel fuer alle Blaetter.
% Einbinden mit % Gemeinsame Praeambel fuer alle Blaetter.
% Einbinden mit % Gemeinsame Praeambel fuer alle Blaetter.
% Einbinden mit % Gemeinsame Praeambel fuer alle Blaetter.
% Einbinden mit \input{../../stil.tex} — Pfad ist relativ zum Kompilierordner.
% Kompiliert wird mit lualatex (wegen luatexja).

\usepackage[ngerman]{babel}
\usepackage[haranoaji]{luatexja-preset}
\usepackage{luatexja-ruby}
\usepackage{booktabs}
\usepackage{array}
\usepackage{xcolor}
\usepackage[most]{tcolorbox}
\usepackage{enumitem}
\usepackage{microtype}
\usepackage{graphicx}
\usepackage{wrapfig}
\usepackage{multicol}
\usepackage{amssymb}
\usepackage{needspace}

% Gedankenstriche als LATEINISCHE Zeichen behandeln.
% luatexja haelt sie sonst fuer Japanisch — und weil japanische Zeichen
% keine Leerzeichen brauchen, frisst es den Zeilenumbruch danach:
% "Nomen —und wenn" statt "Nomen — und wenn".
\ltjdefcharrange{100}{"2013,"2014}
\ltjsetparameter{jacharrange={-100}}
\graphicspath{{../../bilder/}}

\definecolor{akzent}{HTML}{1F4E79}
\definecolor{implizit}{HTML}{1B7A43}
\definecolor{explizit}{HTML}{B3541E}
\definecolor{hell}{HTML}{F4F6F8}

\addtokomafont{disposition}{\color{akzent}}
\setkomafont{title}{\color{akzent}\bfseries}

\newcommand{\impl}{\textcolor{implizit}{\textbf{implizit}}}
\newcommand{\expl}{\textcolor{explizit}{\textbf{explizit}}}
\newcommand{\jterm}[2]{\ruby{#1}{#2}}

% Deutsche Anfuehrungszeichen: babel-Befehle statt Unicode.
% Sonst zwei Probleme: ASCII-" ist bei ngerman ein Umlaut-Shorthand
% ("o -> oe), und luatexja haelt Unicode-Quotes fuer Japanisch und
% setzt Abstand drumrum.
\newcommand{\zitat}[1]{\glqq #1\grqq{}}

% Schreiblinie fuer Antworten
\newcommand{\lin}[1][4cm]{\rule[-0.5em]{#1}{0.4pt}}

\newtcolorbox{merke}[1][]{
  enhanced, colback=hell, colframe=akzent,
  boxrule=0pt, leftrule=3pt, arc=0pt, outer arc=0pt,
  left=8pt, right=8pt, top=6pt, bottom=6pt,
  #1
}
 — Pfad ist relativ zum Kompilierordner.
% Kompiliert wird mit lualatex (wegen luatexja).

\usepackage[ngerman]{babel}
\usepackage[haranoaji]{luatexja-preset}
\usepackage{luatexja-ruby}
\usepackage{booktabs}
\usepackage{array}
\usepackage{xcolor}
\usepackage[most]{tcolorbox}
\usepackage{enumitem}
\usepackage{microtype}
\usepackage{graphicx}
\usepackage{wrapfig}
\usepackage{multicol}
\usepackage{amssymb}
\usepackage{needspace}

% Gedankenstriche als LATEINISCHE Zeichen behandeln.
% luatexja haelt sie sonst fuer Japanisch — und weil japanische Zeichen
% keine Leerzeichen brauchen, frisst es den Zeilenumbruch danach:
% "Nomen —und wenn" statt "Nomen — und wenn".
\ltjdefcharrange{100}{"2013,"2014}
\ltjsetparameter{jacharrange={-100}}
\graphicspath{{../../bilder/}}

\definecolor{akzent}{HTML}{1F4E79}
\definecolor{implizit}{HTML}{1B7A43}
\definecolor{explizit}{HTML}{B3541E}
\definecolor{hell}{HTML}{F4F6F8}

\addtokomafont{disposition}{\color{akzent}}
\setkomafont{title}{\color{akzent}\bfseries}

\newcommand{\impl}{\textcolor{implizit}{\textbf{implizit}}}
\newcommand{\expl}{\textcolor{explizit}{\textbf{explizit}}}
\newcommand{\jterm}[2]{\ruby{#1}{#2}}

% Deutsche Anfuehrungszeichen: babel-Befehle statt Unicode.
% Sonst zwei Probleme: ASCII-" ist bei ngerman ein Umlaut-Shorthand
% ("o -> oe), und luatexja haelt Unicode-Quotes fuer Japanisch und
% setzt Abstand drumrum.
\newcommand{\zitat}[1]{\glqq #1\grqq{}}

% Schreiblinie fuer Antworten
\newcommand{\lin}[1][4cm]{\rule[-0.5em]{#1}{0.4pt}}

\newtcolorbox{merke}[1][]{
  enhanced, colback=hell, colframe=akzent,
  boxrule=0pt, leftrule=3pt, arc=0pt, outer arc=0pt,
  left=8pt, right=8pt, top=6pt, bottom=6pt,
  #1
}
 — Pfad ist relativ zum Kompilierordner.
% Kompiliert wird mit lualatex (wegen luatexja).

\usepackage[ngerman]{babel}
\usepackage[haranoaji]{luatexja-preset}
\usepackage{luatexja-ruby}
\usepackage{booktabs}
\usepackage{array}
\usepackage{xcolor}
\usepackage[most]{tcolorbox}
\usepackage{enumitem}
\usepackage{microtype}
\usepackage{graphicx}
\usepackage{wrapfig}
\usepackage{multicol}
\usepackage{amssymb}
\usepackage{needspace}

% Gedankenstriche als LATEINISCHE Zeichen behandeln.
% luatexja haelt sie sonst fuer Japanisch — und weil japanische Zeichen
% keine Leerzeichen brauchen, frisst es den Zeilenumbruch danach:
% "Nomen —und wenn" statt "Nomen — und wenn".
\ltjdefcharrange{100}{"2013,"2014}
\ltjsetparameter{jacharrange={-100}}
\graphicspath{{../../bilder/}}

\definecolor{akzent}{HTML}{1F4E79}
\definecolor{implizit}{HTML}{1B7A43}
\definecolor{explizit}{HTML}{B3541E}
\definecolor{hell}{HTML}{F4F6F8}

\addtokomafont{disposition}{\color{akzent}}
\setkomafont{title}{\color{akzent}\bfseries}

\newcommand{\impl}{\textcolor{implizit}{\textbf{implizit}}}
\newcommand{\expl}{\textcolor{explizit}{\textbf{explizit}}}
\newcommand{\jterm}[2]{\ruby{#1}{#2}}

% Deutsche Anfuehrungszeichen: babel-Befehle statt Unicode.
% Sonst zwei Probleme: ASCII-" ist bei ngerman ein Umlaut-Shorthand
% ("o -> oe), und luatexja haelt Unicode-Quotes fuer Japanisch und
% setzt Abstand drumrum.
\newcommand{\zitat}[1]{\glqq #1\grqq{}}

% Schreiblinie fuer Antworten
\newcommand{\lin}[1][4cm]{\rule[-0.5em]{#1}{0.4pt}}

\newtcolorbox{merke}[1][]{
  enhanced, colback=hell, colframe=akzent,
  boxrule=0pt, leftrule=3pt, arc=0pt, outer arc=0pt,
  left=8pt, right=8pt, top=6pt, bottom=6pt,
  #1
}
 — Pfad ist relativ zum Kompilierordner.
% Kompiliert wird mit lualatex (wegen luatexja).

\usepackage[ngerman]{babel}
\usepackage[haranoaji]{luatexja-preset}
\usepackage{luatexja-ruby}
\usepackage{booktabs}
\usepackage{array}
\usepackage{xcolor}
\usepackage[most]{tcolorbox}
\usepackage{enumitem}
\usepackage{microtype}
\usepackage{graphicx}
\usepackage{wrapfig}
\usepackage{multicol}
\usepackage{amssymb}
\usepackage{needspace}

% Gedankenstriche als LATEINISCHE Zeichen behandeln.
% luatexja haelt sie sonst fuer Japanisch — und weil japanische Zeichen
% keine Leerzeichen brauchen, frisst es den Zeilenumbruch danach:
% "Nomen —und wenn" statt "Nomen — und wenn".
\ltjdefcharrange{100}{"2013,"2014}
\ltjsetparameter{jacharrange={-100}}
\graphicspath{{../../bilder/}}

\definecolor{akzent}{HTML}{1F4E79}
\definecolor{implizit}{HTML}{1B7A43}
\definecolor{explizit}{HTML}{B3541E}
\definecolor{hell}{HTML}{F4F6F8}

\addtokomafont{disposition}{\color{akzent}}
\setkomafont{title}{\color{akzent}\bfseries}

\newcommand{\impl}{\textcolor{implizit}{\textbf{implizit}}}
\newcommand{\expl}{\textcolor{explizit}{\textbf{explizit}}}
\newcommand{\jterm}[2]{\ruby{#1}{#2}}

% Deutsche Anfuehrungszeichen: babel-Befehle statt Unicode.
% Sonst zwei Probleme: ASCII-" ist bei ngerman ein Umlaut-Shorthand
% ("o -> oe), und luatexja haelt Unicode-Quotes fuer Japanisch und
% setzt Abstand drumrum.
\newcommand{\zitat}[1]{\glqq #1\grqq{}}

% Schreiblinie fuer Antworten
\newcommand{\lin}[1][4cm]{\rule[-0.5em]{#1}{0.4pt}}

\newtcolorbox{merke}[1][]{
  enhanced, colback=hell, colframe=akzent,
  boxrule=0pt, leftrule=3pt, arc=0pt, outer arc=0pt,
  left=8pt, right=8pt, top=6pt, bottom=6pt,
  #1
}
